\chapter{Couche 4 — Les secrets}

La plateforme manipule quatre secrets : la clé de session, la clé du coffre, le
mot de passe de la base, et le jeton des agents. Plus un secret par machine
louée. Aucun ne doit se retrouver en clair, nulle part.

% -----------------------------------------------------------------------------
\mesure{Deux clés, jamais la même}{mesure 8 · \fichier{crypto.py}}

\menace
Une seule clé pour tout est commode : on la perd une fois, on perd tout. Si la
clé qui signe les cookies est aussi celle qui chiffre les mots de passe
\cmd{root}, alors la fuite d'une clé de session --- par nature plus exposée,
puisqu'elle sert à chaque requête --- livre l'accès à toutes les machines
louées.

\parade
Deux clés distinctes, de deux natures différentes, pour deux usages qui n'ont
rien à voir.

\danslecode
\begin{lstlisting}[language=Python]
def load_secret_key() -> str:
    """Cle secrete des sessions Flask."""
    return load_or_create_key("INSACLOUD_SECRET_KEY", ".secret_key",
                              lambda: secrets.token_hex(32))

# Cle de chiffrement des mots de passe root (distincte de la cle de session)
VAULT = Fernet(load_or_create_key("INSACLOUD_VAULT_KEY", ".vault_key",
                                  lambda: Fernet.generate_key().decode()))
\end{lstlisting}

Le commentaire en tête du module le dit explicitement : \emph{« Séparer les
deux garantit qu'un vol de la clé de session ne donne pas accès aux machines
louées. »}

\begin{notion}{D'où viennent les clés}
\cmd{load\_or\_create\_key} suit trois pistes, dans l'ordre :
\begin{enumerate}[leftmargin=16pt,itemsep=2pt]
\item la \textbf{variable d'environnement} --- c'est ainsi qu'Ansible les
      injecte depuis son Vault, en production ;
\item un \textbf{fichier local} déjà présent ;
\item à défaut, une clé \textbf{tirée au hasard}, écrite dans un fichier
      aussitôt mis en \cmd{0600}.
\end{enumerate}
Ce troisième cas est important : il n'existe \emph{aucune} valeur par défaut
inscrite dans le code. Une clé par défaut serait la même sur tous les
déploiements, donc publique.
\end{notion}

\limite
Les fichiers \cmd{.secret\_key} et \cmd{.vault\_key} vivent sur le disque du
contrôleur. Qui lit ce disque les obtient. Un gestionnaire externe
(HashiCorp Vault) les garderait hors de la machine, au prix d'un service
supplémentaire à déployer et à sécuriser lui-même.

% -----------------------------------------------------------------------------
\mesure{Le coffre : Fernet}{mesure 9 · \fichier{crypto.py}, \fichier{services.py}}

\menace
Le mot de passe \cmd{root} de chaque machine louée doit être conservé : pour
l'afficher à l'utilisateur quand il le redemande, et pour le faire tourner. En
clair dans la base, il tombe avec elle.

\parade
Chaque mot de passe est chiffré avant d'être écrit, par \textbf{Fernet} --- la
recette de la bibliothèque \cmd{cryptography} : AES-128 en mode CBC pour le
chiffrement, HMAC-SHA256 pour l'authentification, et un horodatage.

\danslecode
\begin{lstlisting}[language=Python]
def encrypt_secret(value: str) -> str:
    """Chiffre un secret avant de l'ecrire en base."""
    return VAULT.encrypt(value.encode()).decode()
\end{lstlisting}

Appelé aux trois endroits où un mot de passe naît ou change : à la création de
la machine, à la régénération demandée par l'utilisateur, et à la recréation
après la panne d'un nœud.

\begin{notion}{Pourquoi Fernet plutôt qu'AES « à la main »}
Fernet est un format \textbf{authentifié} : il ne se contente pas de chiffrer,
il signe. Un attaquant qui modifierait un octet du texte chiffré en base
provoquerait un \cmd{InvalidToken} au déchiffrement, et non un mot de passe
différent mais accepté. Chiffrer sans authentifier est une erreur classique,
qui ouvre des attaques par altération du texte chiffré ; Fernet la rend
impossible par construction.
\end{notion}

\limite
Le coffre protège la base, pas la mémoire : pendant qu'il sert une page, le
processus détient le mot de passe en clair. Et la clé du coffre est sur la même
machine que la base --- la séparation est logique, pas physique.

% -----------------------------------------------------------------------------
\mesure{La fenêtre du coffre}{mesure 10 · \fichier{app.py}}

\menace
Afficher en permanence les mots de passe \cmd{root} sur le tableau de bord,
c'est les offrir à quiconque passe derrière l'utilisateur --- ou à une session
restée ouverte.

\parade
Les mots de passe sont \textbf{masqués par défaut}. Pour les voir, l'utilisateur
doit ressaisir le mot de passe de \emph{son compte} ; le coffre s'ouvre alors
pour cinq minutes, puis se referme seul.

\danslecode
\begin{lstlisting}[language=Python]
VAULT_WINDOW_MINUTES = int(os.environ.get("INSACLOUD_VAULT_WINDOW", "5"))
\end{lstlisting}

Le déchiffrement n'a même pas lieu si la fenêtre est fermée : le gabarit ne
reçoit pas le secret, il reçoit \cmd{None}. Un secret qui n'est pas envoyé au
navigateur ne peut pas y être intercepté.

\begin{lstlisting}[language=Python]
secret = decrypt_secret(row["root_password"]) if (vault_open and running) else None
\end{lstlisting}

\limite
Cinq minutes suffisent à une capture d'écran. La fenêtre réduit l'exposition,
elle ne la supprime pas --- et c'est précisément pourquoi l'authentification
par clé publique (couche 5) est le meilleur chemin : elle rend le mot de passe
inutile.

% -----------------------------------------------------------------------------
\mesure{Ansible Vault}{mesure 11 · \fichier{3\_ansible/group\_vars/}}

\menace
Les secrets doivent arriver sur le serveur. Trois mauvaises façons de s'y
prendre : les écrire en clair dans le dépôt, les taper à la main à chaque
déploiement, ou les passer en variables d'environnement depuis un terminal ---
où ils restent dans l'historique du shell.

\parade
Un fichier chiffré en \textbf{AES-256}, versionné tel quel dans le dépôt,
qu'Ansible déchiffre au moment de jouer le playbook avec une phrase de passe
fournie à l'exécution.

\danslecode
\fichier{3\_ansible/group\_vars/insacloud/vault.yml} contient la clé de session,
la clé du coffre et le mot de passe de la base. Le rôle \cmd{webapp} vérifie
même leur \emph{forme} avant de déployer --- une clé Fernet doit faire
exactement 44 caractères --- plutôt que de laisser l'application échouer au
démarrage avec un message obscur.

La phrase de passe vit dans un fichier \cmd{.vault\_pass} en \cmd{0600}, exclu
du dépôt par le \fichier{.gitignore}.

\begin{retenir}
Le fichier chiffré est \emph{dans} le dépôt, et c'est l'intérêt : les secrets
voyagent avec le code, le déploiement reste entièrement automatisable, et il
n'y a aucune étape manuelle de copie de clés --- étape qui, en pratique, finit
toujours par être faite de travers.
\end{retenir}

\limite
La phrase de passe est le dernier maillon, et elle n'est protégée par rien
d'autre que les droits du fichier. Qui l'obtient obtient tout le reste.

% -----------------------------------------------------------------------------
\mesure{Les secrets sur le serveur}{mesure 12 · rôle \cmd{webapp}}

\menace
Un secret déposé sur le serveur peut fuir par trois chemins qu'on oublie :
l'unité systemd (lisible par tous), les journaux, et la liste des processus.

\parade
Les secrets sont déposés dans un \textbf{fichier séparé en \cmd{0600}}, lu par
systemd via \cmd{EnvironmentFile} --- jamais écrits dans l'unité elle-même, qui
est en lecture publique, et jamais passés en arguments de ligne de commande,
qui apparaissent dans \cmd{ps}.

\danslecode
Du côté des machines louées, le même soin : l'\fichier{entrypoint.sh} termine
par

\begin{lstlisting}[language=bash]
# Le mot de passe ne doit pas rester lisible dans l'environnement des processus
unset ROOT_PASSWORD SSH_PUBKEY
exec "$(command -v supervisord)" -c "$CONF"
\end{lstlisting}

Sans ce \cmd{unset}, le mot de passe \cmd{root} resterait lisible dans
\cmd{/proc/1/environ} pendant toute la vie de la machine.

\limite
\cmd{docker inspect} montre les variables d'environnement passées au
\cmd{docker run}, donc le mot de passe initial, à qui a accès au démon Docker.
Mais qui a accès au démon Docker a déjà tout (chapitre 1) : ce n'est pas une
fuite supplémentaire.

% -----------------------------------------------------------------------------
\mesure{La recherche de secrets commis}{mesure 13 · \cmd{gitleaks}}

\menace
Un secret commis par erreur puis retiré au commit suivant \textbf{reste dans
l'historique}, consultable par quiconque clone le dépôt. Le retirer de la
dernière version ne sert à rien.

\parade
\cmd{gitleaks} parcourt \textbf{tout l'historique} à chaque envoi, dans un
travail dédié de la chaîne d'intégration.

\danslecode
Le travail \cmd{secrets} du fichier \fichier{.github/workflows/ci.yml}. Le
\fichier{.gitignore} tient par ailleurs hors du dépôt les clés
(\cmd{.secret\_key}, \cmd{.vault\_key}), le fichier \cmd{.env} de la pile
compose, et la phrase de passe du Vault.

\limite
\cmd{gitleaks} reconnaît les formats connus --- clés AWS, jetons GitHub,
clés privées. Un secret de forme quelconque, comme un mot de passe de base de
données, peut lui échapper. D'où la règle, humaine : aucun secret n'entre dans
un fichier versionné, chiffré ou non.
